\documentclass[12pt,a4paper]{report}

% Praktikumsdokumentation MV, Modul 1 "Delivery".
% Layout angelehnt an die TH-OWL-Formatvorlage (Arial/Helvetica 12 pt).
% Gliederung gemäß Inhalt_Dokumentation.txt.

% Sprache, Schrift und Mikrotypografie
\usepackage{iftex}
\ifPDFTeX
  \usepackage[utf8]{inputenc}
  \usepackage[T1]{fontenc}
  \usepackage[scaled=0.95]{helvet}
  \renewcommand{\familydefault}{\sfdefault}
\else
  \usepackage{fontspec}
  \IfFontExistsTF{Arial}{%
    \setmainfont{Arial}
    \setsansfont{Arial}
  }{%
    \setmainfont{TeX Gyre Heros}
    \setsansfont{TeX Gyre Heros}
  }
\fi
% shorthands=off: sonst frisst babel das " in JSON-Beispielen wie {"value": true}
\usepackage[ngerman,shorthands=off]{babel}
\usepackage{microtype}
\usepackage{setspace}

% Seitenlayout: A4, oben 2,5 cm, unten 2 cm, rechts 2 cm, links 2 cm + 1 cm Bundsteg
\usepackage[
  a4paper,
  top=2.5cm,
  bottom=2cm,
  left=3cm,
  right=2cm,
  headheight=16pt,
  headsep=0.8cm,
  footskip=1.2cm
]{geometry}

% Tabellen, Grafiken, Verweise
\usepackage{graphicx}
\usepackage{booktabs}
\usepackage{array}
\usepackage{tabularx}
\usepackage{longtable}
\usepackage{enumitem}
\usepackage{caption}
\usepackage{float}
\usepackage{fancyhdr}
\usepackage{titlesec}
\usepackage{tocloft}
\usepackage[table]{xcolor}
\usepackage{listings}
\usepackage[pdfpagelabels,plainpages=false,hypertexnames=false,hidelinks]{hyperref}

% Diagramme: Zustandsautomat, Topologie, Zeitverläufe
\usepackage{tikz}
\usetikzlibrary{automata,positioning,arrows.meta,shapes.geometric,fit,calc}
\usepackage{tikz-timing}

% Grundformat: Blocksatz, mind. 16 pt Zeilenabstand, 8 pt nach Absätzen
\setstretch{1.104}
\setlength{\parindent}{0pt}
\setlength{\parskip}{8pt}
\frenchspacing
\sloppy

% Gliederung: dekadisch, ohne Punkt am Ende
\setcounter{secnumdepth}{3}
\setcounter{tocdepth}{2}
\titleformat{\chapter}[hang]
  {\normalfont\linespread{1}\fontsize{14pt}{16pt}\selectfont\bfseries}
  {\thechapter}{1em}{}
\titleformat{\section}
  {\normalfont\linespread{1}\fontsize{12pt}{16pt}\selectfont\bfseries}
  {\thesection}{1em}{}
\titleformat{\subsection}
  {\normalfont\linespread{1}\fontsize{12pt}{16pt}\selectfont}
  {\thesubsection}{1em}{}
\titlespacing*{\chapter}{0pt}{18pt}{9pt}
\titlespacing*{\section}{0pt}{18pt}{9pt}
\titlespacing*{\subsection}{0pt}{18pt}{6pt}

% Verzeichnisse
\renewcommand{\contentsname}{Inhaltsverzeichnis}
\renewcommand{\listfigurename}{Abbildungsverzeichnis}
\renewcommand{\listtablename}{Tabellenverzeichnis}
\renewcommand{\bibname}{Literaturverzeichnis}
\renewcommand{\cftsecleader}{\cftdotfill{\cftdotsep}}
\renewcommand{\cfttoctitlefont}{\normalfont\linespread{1}\fontsize{14pt}{16pt}\selectfont\bfseries}
\renewcommand{\cftloftitlefont}{\normalfont\linespread{1}\fontsize{14pt}{16pt}\selectfont\bfseries}
\renewcommand{\cftlottitlefont}{\normalfont\linespread{1}\fontsize{14pt}{16pt}\selectfont\bfseries}

% Beschriftungen: Nummerierung je Kapitel
\counterwithin{figure}{chapter}
\counterwithin{table}{chapter}
\captionsetup{
  font=normalsize,
  labelfont=bf,
  labelsep=colon,
  justification=raggedright,
  singlelinecheck=false
}

% Aufzählungen: linksbündig, hängend eingerückt
\setlist[itemize]{label=\textbullet, leftmargin=1.25cm, labelsep=0.35cm, itemsep=0pt, parsep=0pt, topsep=0pt}
\setlist[enumerate]{leftmargin=1.25cm, labelsep=0.35cm, itemsep=0pt, parsep=0pt, topsep=0pt}

% Quelltext und JSON
\lstdefinelanguage{json}{
  morestring=[b]",
  literate={:}{{\textcolor{gray}{:}}}1 {,}{{\textcolor{gray}{,}}}1,
}
\lstset{
  basicstyle=\ttfamily\small,
  columns=fullflexible,
  keepspaces=true,
  frame=single,
  rulecolor=\color{gray!60},
  backgroundcolor=\color{gray!8},
  xleftmargin=0.2cm,
  framexleftmargin=0.2cm,
  breaklines=true,
  stringstyle=\color{teal!70!black},
  showstringspaces=false
}

% ---------------------------------------------------------------------------
% Metadaten (hier anpassen)
% ---------------------------------------------------------------------------
\newcommand{\Hochschule}{Technische Hochschule Ostwestfalen-Lippe}
\newcommand{\HochschuleKurz}{TH OWL}
\newcommand{\Fachbereich}{Fachbereich Informatik und Automation}
\newcommand{\Lehrveranstaltung}{Maschinennahe Vernetzung - Praktikum}
\newcommand{\Arbeitstyp}{Praktikumsdokumentation}
\newcommand{\Arbeitstitel}{Steuerung von Modul~1 \glqq Delivery\grqq{} der Popcorn-Modellanlage über MQTT}
\newcommand{\Semester}{WS 2026/27}
\newcommand{\Abgabedatum}{\today}
\newcommand{\AutorNamenFuss}{Braun, Mahnke, Porde}

% Bearbeiter mit Matrikelnummer
\newcommand{\AutorenTabelle}{%
  \begin{tabular}{@{}ll@{}}
    Jeremy Braun & Matr.-Nr.: 2012 6818 \\
    Jan Mahnke   & Matr.-Nr.: 1546 7074 \\
    Dustin Porde & Matr.-Nr.: 2013 4105 \\
  \end{tabular}%
}

% Betreuer
\newcommand{\BetreuerTabelle}{%
  \begin{tabular}{@{}ll@{}}
    Prof. Dr. Jürgen Jasperneite \\
    Benedikt Lücke \\
  \end{tabular}%
}

% ---------------------------------------------------------------------------
% Kopf- und Fußzeile
% ---------------------------------------------------------------------------
\pagestyle{fancy}
\fancyhf{}
\renewcommand{\headrulewidth}{0pt}
\renewcommand{\footrulewidth}{0pt}
\renewcommand{\chaptermark}[1]{\markboth{#1}{}}
\fancyhead[L]{\nouppercase{\leftmark}}
\fancyhead[R]{\thepage}
\fancyfoot[L]{\AutorNamenFuss}
\fancyfoot[C]{\HochschuleKurz}
\fancyfoot[R]{\Semester}
\fancypagestyle{plain}{%
  \fancyhf{}%
  \fancyhead[L]{\nouppercase{\leftmark}}\fancyhead[R]{\thepage}%
  \fancyfoot[L]{\AutorNamenFuss}\fancyfoot[C]{\HochschuleKurz}\fancyfoot[R]{\Semester}%
}

% ---------------------------------------------------------------------------
% Hilfsbefehle
% ---------------------------------------------------------------------------

% Deckblatt
\newcommand{\makedeckblatt}{%
  \begin{titlepage}
    \thispagestyle{empty}
    \centering
    \vspace*{1.0cm}
    {\fontsize{23pt}{28pt}\selectfont\bfseries \Hochschule\par}
    \vspace{0.6cm}
    {\fontsize{16pt}{20pt}\selectfont \Fachbereich\par}
    \vspace{1.6cm}
    {\fontsize{22pt}{26pt}\selectfont\bfseries \Arbeitstyp\par}
    \vspace{0.4cm}
    {\fontsize{14pt}{18pt}\selectfont \Lehrveranstaltung\par}
    \vspace{1.4cm}
    {\fontsize{18pt}{24pt}\selectfont\bfseries \Arbeitstitel\par}
    \vfill
    \begin{flushleft}
      \fontsize{13pt}{18pt}\selectfont
      \begin{tabular}{@{}p{3.4cm}@{\hspace{0.4cm}}l@{}}
        \textbf{Bearbeiter:} & \AutorenTabelle \\[1cm]
        \textbf{Betreuer:}   & \BetreuerTabelle \\[0.6cm]
        \textbf{Semester:}   & \Semester \\[0.3cm]
        \textbf{Abgabe:}     & \Abgabedatum \\
      \end{tabular}
    \end{flushleft}
    \vspace{1cm}
  \end{titlepage}
}

% Ungezähltes Kapitel im Vorspann
\newcommand{\vorspannkapitel}[1]{%
  \clearpage
  \phantomsection
  \chapter*{#1}
  \markboth{#1}{}
}

% Abkürzungen
\newenvironment{abkuerzungsverzeichnis}{%
  \vorspannkapitel{Abkürzungsverzeichnis}
  \addcontentsline{toc}{chapter}{Abkürzungsverzeichnis}
  \begin{longtable}{@{}p{3cm}@{\hspace{0.5cm}}p{\dimexpr\textwidth-3.5cm\relax}@{}}
}{%
  \end{longtable}
}
\newcommand{\abk}[2]{\leavevmode\hypertarget{abk:#1}{#1} & #2\\}
\DeclareRobustCommand{\abkref}[1]{\hyperlink{abk:#1}{#1}}

% Bild einbinden oder Platzhalter anzeigen, solange die Datei fehlt
% \bildoderplatzhalter[Breite]{abbildungen/datei}{Beschriftung}{label}
\newcommand{\bildoderplatzhalter}[4][width=\linewidth]{%
  \begin{figure}[htbp]
    \centering
    \IfFileExists{#2}{%
      \includegraphics[#1]{#2}%
    }{%
      \fbox{\parbox[c][4cm][c]{0.85\linewidth}{\centering Abbildung fehlt:\\ \texttt{\detokenize{#2}}}}%
    }%
    \caption{#3}
    \label{#4}
  \end{figure}
}

% Lampe als Symbol: \lampe{red}{1} = an, \lampe{red}{0} = aus
\newcommand{\lampe}[2]{%
  \ifnum#2=1
    \tikz[baseline=-0.6ex]\filldraw[draw=black, fill=#1] (0,0) circle (0.17cm);%
  \else
    \tikz[baseline=-0.6ex]\draw[gray!70] (0,0) circle (0.17cm);%
  \fi
}
\definecolor{lamperot}{RGB}{220,30,30}
\definecolor{lampeorange}{RGB}{245,150,20}
\definecolor{lampegruen}{RGB}{40,170,60}

% Zeitdiagramme (tikz-timing): einheitlicher Rahmen
% \begin{zeitdiagramm}{Beschriftung}{label} ... Zeilen ... \end{zeitdiagramm}
\tikzset{timing/d/text/.style={font=\scriptsize}}
\newenvironment{zeitdiagramm}[2]{%
  \def\zeitdiagrammcaption{#1}%
  \def\zeitdiagrammlabel{#2}%
  % [H]: Diagramm steht immer direkt unter dem zugehörigen Text
  \begin{figure}[H]
  \centering
  \small
}{%
  \caption{\zeitdiagrammcaption}
  \label{\zeitdiagrammlabel}
  \end{figure}
}

% TODO-Markierung im Text (vor Abgabe alle entfernen)
\newcommand{\todo}[1]{\textcolor{red}{\textbf{TODO:} #1}}

% ---------------------------------------------------------------------------
\begin{document}

\makedeckblatt

\pagenumbering{Roman}
\setcounter{page}{1}

\clearpage
\phantomsection
\markboth{Inhaltsverzeichnis}{}
\tableofcontents

\clearpage
\phantomsection
\markboth{Abbildungsverzeichnis}{}
\addcontentsline{toc}{chapter}{Abbildungsverzeichnis}
\listoffigures

\clearpage
\phantomsection
\markboth{Tabellenverzeichnis}{}
\addcontentsline{toc}{chapter}{Tabellenverzeichnis}
\listoftables

\begin{abkuerzungsverzeichnis}
  \abk{IP}{\textbf{I}nternet \textbf{P}rotocol}
  \abk{JSON}{\textbf{J}ava\textbf{S}cript \textbf{O}bject \textbf{N}otation}
  \abk{MQTT}{\textbf{M}essage \textbf{Q}ueuing \textbf{T}elemetry \textbf{T}ransport}
  \abk{NC}{\textbf{N}ormally \textbf{C}losed, Öffner}
  \abk{NO}{\textbf{N}ormally \textbf{O}pen, Schließer}
  \abk{PAP}{\textbf{P}rogramm\textbf{a}blauf\textbf{p}lan}
  \abk{PWM}{\textbf{P}uls\textbf{w}eiten\textbf{m}odulation}
  \abk{QoS}{\textbf{Q}uality \textbf{o}f \textbf{S}ervice}
  \abk{TCP}{\textbf{T}ransmission \textbf{C}ontrol \textbf{P}rotocol}
  \abk{VPS}{Modellanlage des Praktikums (Popcorn-Anlage) \todo{Langform klären}}
\end{abkuerzungsverzeichnis}


\clearpage
\pagenumbering{arabic}
\setcounter{page}{1}

\chapter{Einleitung}
\label{chap:einleitung}

% Inhalt laut Vorgabe:
%  - Beschreibung der gesamten Anlage (Kurzfassung)
%  - Beschreibung des eigenen Moduls (detailliert)

\section{Gesamtanlage}
\label{sec:gesamtanlage}

Die Modellanlage (\abkref{VPS}) bildet die Herstellung von Popcorn in vier
hintereinander geschalteten Modulen ab. Der Mais durchläuft die Module in der
Reihenfolge Delivery, Storage, Dosing und Production. Die Anlage wird
ausschließlich über einen \abkref{MQTT}-Broker angesprochen: Sensorwerte werden
vom Broker gelesen, Aktoren werden durch Nachrichten an den Broker geschaltet
\cite{doku-popcorn}.

\begin{table}[htbp]
  \centering
  \begin{tabularx}{\textwidth}{@{}l l X@{}}
    \toprule
    \textbf{Modul} & \textbf{Name} & \textbf{Aufgabe} \\
    \midrule
    1 & Delivery   & Anlieferung: Mais aus dem Silo per Förderband und Druckluft an Modul~2 übergeben \\
    2 & Storage    & Zwischenlagerung; Weitertransport über pneumatischen Muskel und Druckluft \\
    3 & Dosing     & Dosierung über Förderschnecke (Drehzahl per \abkref{PWM}) und Wägezelle \\
    4 & Production & Poppen mit Heizspirale; Becher und Flanke für die Ausgabe \\
    \bottomrule
  \end{tabularx}
  \caption{Module der Gesamtanlage}
  \label{tab:module}
\end{table}

\todo{Kurzbeschreibung der Gesamtanlage ergänzen: Zusammenspiel der Module,
welche Gruppe welches Modul steuert, Material- und Signalfluss.}

\section{Modul 1 -- Delivery}
\label{sec:modul-delivery}

Modul~1 ist das erste Modul der Anlage. Es fördert den Mais aus seinem Silo über
ein Förderband in einen Trichter. Von dort wird der Mais mit Druckluft durch einen
Schlauch zu Modul~2 (Storage) transportiert. Sobald der Motor des Förderbands
eingeschaltet wird, läuft die Druckluft automatisch mit \cite{doku-popcorn}.

Das Modul besitzt fünf Aktoren (Motor, Druckluft und drei Signallampen) sowie vier
Sensoren (Silo oben, Silo unten, Trichter und Not-Aus). Eine vollständige Liste mit
Funktion und MQTT-Topic steht in Kapitel~\ref{chap:hardware}.

\bildoderplatzhalter[width=0.8\linewidth]{abbildungen/delivery-schema.png}
  {Aufbau von Modul~1 (Delivery) mit Lage der Sensoren und Aktoren \cite{doku-popcorn}}
  {fig:delivery-schema}

Die Steuerung ist in Python mit der Bibliothek \texttt{paho-mqtt} umgesetzt
\cite{paho}. Sie besteht aus drei Schichten:
\begin{itemize}
  \item \textbf{Komponenten} (\texttt{komponenten.py}): kapseln je einen Aktor
        und senden dessen Befehle an den Broker.
  \item \textbf{Modul} (\texttt{Module/delivery.py}): hält die aktuellen
        Sensorwerte und verteilt eingehende Nachrichten.
  \item \textbf{Zustandsautomat} (\texttt{zustandsautomat.py}): entscheidet
        zyklisch alle 0,1~s anhand der Sensorwerte über den nächsten
        Zustand und schaltet die Aktoren über das Modul.
\end{itemize}

\todo{Aufgabe und Randbedingungen des Moduls detailliert beschreiben
(Startbedingung, Abschaltbedingungen, Kopplung zu Modul~2, Sicherheitsverhalten).}

\chapter{Programmablauf}
\label{chap:programmablauf}

% Inhalt laut Vorgabe:
%  - PAP oder Zustandsdiagramm o. Ä.
%  - ggf. Erläuterungen zum Bild

Die Steuerung von Modul~1 ist als Zustandsautomat umgesetzt. Die Methode
\texttt{schritt()} wird alle 0,1~s aufgerufen. Sie prüft zuerst die
Sicherheitsbedingungen und wertet danach den aktuellen Zustand aus.
Abbildung~\ref{fig:zustandsautomat} zeigt die Zustände und Übergänge.

\begin{figure}[htbp]
  \centering
  \begin{tikzpicture}[
      zustand/.style={draw, rectangle, rounded corners=4pt, minimum width=2.6cm,
                      minimum height=1.1cm, align=center, font=\small\bfseries},
      stopzustand/.style={zustand, draw=lamperot, very thick},
      uebergang/.style={-{Stealth[length=2.5mm]}, thick},
      bed/.style={font=\scriptsize, align=center},
      node distance=2.2cm and 3.2cm
    ]
    \node[zustand]                         (init)     {INIT};
    \node[zustand, right=of init]          (bereit)   {BEREIT};
    \node[zustand, right=of bereit]        (vorlauf)  {VORLAUF};
    \node[zustand, below=of vorlauf]       (laeuft)   {LÄUFT};
    \node[zustand, below=of bereit]        (nachlauf) {NACHLAUF};
    \node[stopzustand, below=of init]      (stop)     {STOP};

    \draw[uebergang] ($(init.west)+(-0.8,0)$) -- (init.west);

    \draw[uebergang] (init) -- node[bed, above] {Aktoren aus} (bereit);
    \draw[uebergang] (bereit) -- node[bed, above] {Silo nicht leer $\wedge$\\ M2 nicht voll}
                                node[bed, below] {/ Luft an} (vorlauf);
    \draw[uebergang] (vorlauf) -- node[bed, right] {nach 1~s\\ / Motor an} (laeuft);
    \draw[uebergang] (laeuft) -- node[bed, below] {Silo leer $\vee$\\ M2 voll} (nachlauf);
    \draw[uebergang] (nachlauf) -- node[bed, right] {nach 1~s\\ / Motor, Luft aus} (bereit);
    \draw[uebergang] (stop) -- node[bed, above, sloped] {kein Sicherheitsfehler} (bereit);

    \node[bed, below=0.9cm of stop, text width=5.2cm, draw=lamperot, dashed,
          rounded corners=3pt, inner sep=4pt] (alle)
      {aus jedem Zustand:\\ Not-Aus $\vee$ M2-Not-Aus $\vee$ Trichter voll\\ / Motor, Luft aus};
    \draw[uebergang, draw=lamperot] (alle) -- (stop);
  \end{tikzpicture}
  \caption{Zustandsautomat von Modul~1 (Delivery)}
  \label{fig:zustandsautomat}
\end{figure}

\section{Zustände}

\begin{table}[htbp]
  \centering
  \begin{tabularx}{\textwidth}{@{}l X@{}}
    \toprule
    \textbf{Zustand} & \textbf{Beschreibung} \\
    \midrule
    INIT     & Startzustand nach dem Einschalten. Alle Aktoren werden ausgeschaltet, danach folgt BEREIT. \\
    BEREIT   & Wartet, bis im Silo Material liegt und Modul~2 Material aufnehmen kann. Dann wird die Druckluft eingeschaltet. \\
    VORLAUF  & Die Druckluft läuft 1~s vor, damit der Schlauch frei ist, bevor das Förderband Mais nachschiebt. \\
    LÄUFT    & Förderband und Druckluft laufen. Ist das Silo leer oder Modul~2 voll, folgt der NACHLAUF. \\
    NACHLAUF & Förderband und Druckluft laufen 1~s nach, danach werden beide abgeschaltet. \\
    STOP     & Sicherheitszustand. Alle Aktoren sind aus, bis kein Sicherheitsfehler mehr anliegt. \\
    \bottomrule
  \end{tabularx}
  \caption{Zustände des Automaten}
  \label{tab:zustaende}
\end{table}

\section{Sicherheitsbedingungen}

Vor jeder Zustandsauswertung prüft der Automat drei Sicherheitsbedingungen:
den eigenen Not-Aus, den Not-Aus von Modul~2 und einen vollen Trichter. Trifft
eine davon zu, werden Motor und Druckluft im selben Schritt ausgeschaltet und
der Automat wechselt nach STOP, unabhängig vom aktuellen Zustand.


\chapter{Hardware und Topologie}
\label{chap:hardware}

% Inhalt laut Vorgabe:
%  - Sensor-/Aktorliste mit Funktion des Sensors/Aktors
%  - Topologie des gesamten Systems (inkl. genutzter Protokolle)

\section{Sensoren und Aktoren von Modul 1}
\label{sec:sensoren-aktoren}

Alle Topics und Standardwerte stammen aus der VPS Simple View
\cite{vps-simple-view}. Sensoren senden ihren Wert als \abkref{JSON}-Objekt
\texttt{\{"value":true,"dataType":"Boolean"\}}, Motor und Druckluft werden mit dem Text
\texttt{true}/\texttt{false} geschaltet (Details in Kapitel~\ref{chap:kommunikation}).

\begin{table}[htbp]
  \centering
  \footnotesize
  \begin{tabularx}{\textwidth}{@{}l l l X@{}}
    \toprule
    \textbf{Sensor} & \textbf{Topic} & \textbf{Default} & \textbf{Funktion / Bedeutung von \texttt{true}} \\
    \midrule
    Silo Top     & \texttt{.../Top}       & \texttt{true}  & Füllstand oben: Silo ist voll \\
    Silo Bottom  & \texttt{.../Bottom}    & \texttt{true}  & Füllstand unten: Material vorhanden, Silo nicht leer \\
    Trichter     & \texttt{.../Funnel}    & \texttt{false} & Trichter am Schlauch ist voll, Förderung muss stoppen \\
    Not-Aus      & \texttt{.../Emergency} & \texttt{true}  & Not-Aus-Taster ist \emph{nicht} gedrückt (Normalbetrieb) \\
    \bottomrule
  \end{tabularx}
  \caption[Sensoren von Modul~1]{Sensoren von Modul~1, Topic-Präfix \texttt{MV/VPS/Sensors/Delivery/}}
  \label{tab:sensoren}
\end{table}

Der Not-Aus ist aus Sicherheitsgründen invertiert ausgelegt: Im Normalbetrieb
meldet er \texttt{true}, erst \texttt{false} bedeutet \glqq Not-Aus gedrückt\grqq.
Nach dem Ruhestromprinzip führt so auch ein Drahtbruch am Taster zu \texttt{false}
und damit zum Stopp.

\begin{table}[htbp]
  \centering
  \footnotesize
  \begin{tabularx}{\textwidth}{@{}l l l X@{}}
    \toprule
    \textbf{Aktor} & \textbf{Topic} & \textbf{Typ} & \textbf{Funktion} \\
    \midrule
    Förderband (Motor) & \texttt{.../Transport/Actor} & bool & Fördert Mais aus dem Silo in den Trichter. Schaltet die Druckluft automatisch mit ein. \\
    Druckluft          & \texttt{.../Transport/Air}   & bool & Bläst den Mais durch den Schlauch zu Modul~2 \\
    Lampen             & \texttt{.../Lamps}           & JSON & Signalsäule rot/orange/grün zur Zustandsanzeige (Kapitel~\ref{chap:signale}) \\
    \bottomrule
  \end{tabularx}
  \caption[Aktoren von Modul~1]{Aktoren von Modul~1, Topic-Präfix \texttt{MV/VPS/Delivery/}}
  \label{tab:aktoren}
\end{table}

% Neue Seite: Text, Abbildung und Protokollliste von 3.2 stehen zusammen
\clearpage
\section{Topologie}
\label{sec:topologie}

Abbildung~\ref{fig:topologie} zeigt die Kommunikationswege. Die Python-Steuerung
und die Modellanlage tauschen keine Daten direkt aus. Beide sind Clients des
MQTT-Brokers, der die Nachrichten anhand der Topics weiterleitet
\cite{mqtt311}.

% [H]: Abbildung direkt unter dem Text, nicht allein auf einer Float-Seite
\begin{figure}[H]
  \centering
  \begin{tikzpicture}[
      kasten/.style={draw, rounded corners=4pt, align=center, font=\small,
                     minimum height=1.0cm, inner sep=5pt},
      gruppe/.style={draw, dashed, rounded corners=6pt, inner sep=8pt},
      verbindung/.style={{Stealth[length=2.2mm]}-{Stealth[length=2.2mm]}, thick},
      prot/.style={font=\scriptsize, align=center},
      node distance=0.5cm
    ]
    % Steuerungs-PC
    \node[kasten, minimum width=3.4cm] (m1) {\texttt{delivery.py}\\ Modul 1};
    \node[kasten, minimum width=3.4cm, below=of m1] (m2) {\texttt{storage.py}\\ Modul 2};
    \node[gruppe, fit=(m1)(m2), label={[font=\small\bfseries]above:Steuerungs-PC}] (pc) {};
    \node[font=\scriptsize, below=0.1cm of pc] {Python 3.12, paho-mqtt 2.1};

    % Broker
    \node[kasten, fill=gray!10, minimum width=3.6cm, minimum height=1.6cm, right=2.0cm of pc]
      (broker) {\textbf{MQTT-Broker}\\ \texttt{dover.init.th-owl.de}\\ Port 1884};

    % Anlage
    \node[kasten, minimum width=3.0cm, right=2.0cm of broker, yshift=1.35cm] (vm1) {M1 Delivery};
    \node[kasten, minimum width=3.0cm, below=0.25cm of vm1] (vm2) {M2 Storage};
    \node[kasten, minimum width=3.0cm, below=0.25cm of vm2] (vm3) {M3 Dosing};
    \node[kasten, minimum width=3.0cm, below=0.25cm of vm3] (vm4) {M4 Production};
    \node[gruppe, fit=(vm1)(vm4), label={[font=\small\bfseries]above:Modellanlage (VPS)}] (vps) {};

    \draw[verbindung] (pc) -- node[prot, above] {MQTT 3.1.1}
                             node[prot, below] {TCP/IP} (broker);
    \draw[verbindung] (broker) -- node[prot, above] {MQTT}
                                 node[prot, below] {TCP/IP} (vps);

    % Weitere Gruppen
    \node[kasten, dashed, below=1.3cm of broker, minimum width=3.6cm] (andere)
      {Steuerungen M3/M4\\ (andere Gruppen)};
    \draw[verbindung, dashed] (andere) -- (broker);
  \end{tikzpicture}
  \caption{Topologie des Gesamtsystems mit den genutzten Protokollen}
  \label{fig:topologie}
\end{figure}

\begin{itemize}
  \item \textbf{Transport:} \abkref{TCP}/\abkref{IP} über das Hochschulnetz,
        Broker-Port 1884.
  \item \textbf{Anwendungsprotokoll:} MQTT 3.1.1 mit Benutzername und Passwort,
        \abkref{QoS}~0, Publish/Subscribe über Topics.
  \item \textbf{Nutzdaten:} Text \texttt{true}/\texttt{false} bzw. \abkref{JSON}
        \cite{rfc8259}.
\end{itemize}


\chapter{Zeitlicher Verlauf der Sensoren und Aktoren}
\label{chap:zeitverlauf}

% Inhalt laut Vorgabe:
%  - Diagramm mit allen Sensoren und Aktoren über der Zeit
%  - 1. normaler Ablauf, mindestens 3 Fehlerfälle
%
% Hinweise zu tikz-timing:
%   H = high (true), L = low (false), D{Text} = Bus/Zustand mit Beschriftung
%   Zahl davor = Länge in Rastereinheiten, z. B. 3H = drei Einheiten high.
%   Jede Zeile muss in Summe gleich lang sein (hier 14 Einheiten = 7 s).

Die folgenden Diagramme zeigen alle Sensoren und Aktoren von Modul~1 sowie die
Handshake-Signale von Modul~2 über der Zeit. Ein Rasterkästchen entspricht
0,5~s. Hoch bedeutet \texttt{true}: bei Silo- und Trichtersensoren
\glqq erkannt\grqq, bei Aktoren \glqq an\grqq. Die beiden Not-Aus-Signale sind
sicherheitsgerichtet ausgelegt (Kapitel~\ref{sec:sensoren-aktoren}): Hoch
bedeutet \glqq Not-Aus nicht gedrückt\grqq, erst ein Wechsel auf low löst den
Stopp aus.

INIT und der direkt folgende Zustand BEREIT (\glqq BER.\grqq) dauern in der
Steuerung nur je einen Zyklus von 0,1~s. Sie sind zur besseren Lesbarkeit
mit einem Rasterkästchen dargestellt.

\section{Normaler Ablauf}
\label{sec:verlauf-normal}

Material liegt im Silo, Modul~2 ist frei. Die Steuerung startet die Druckluft,
schaltet nach 1~s das Förderband zu und fördert. Während der Förderung wird das
Silo kurzzeitig ganz gefüllt (Silo Top), anschließend leert es sich, bis auch
Silo Bottom kein Material mehr erkennt. Danach laufen Band und Luft 1~s nach.
Silo Top hat dabei keinen Einfluss auf den Ablauf.

Hinweis: Die Förderphase ist stark verkürzt dargestellt. In der realen Anlage
leert sich ein volles Silo nicht innerhalb von 2,5~s.

\begin{zeitdiagramm}{Zeitverlauf im normalen Ablauf}{fig:verlauf-normal}
  \resizebox{\textwidth}{!}{%
  \begin{tikztimingtable}[timing/slope=0, xscale=2.6, yscale=1.4, font=\small]
    Zustand        & 1D{INIT} 1D{BER.} 2D{VORLAUF} 5D{LÄUFT} 2D{NACHLAUF} 3D{BEREIT} \\
    Silo Bottom    & 9H 5L \\
    Silo Top       & 5L 2H 7L \\
    Trichter       & 14L \\
    Not-Aus        & 14H \\
    M2 Silo voll   & 14L \\
    M2 Not-Aus     & 14H \\
    Motor          & 4L 7H 3L \\
    Druckluft      & 2L 9H 3L \\
    Lampe rot      & 1H 13L \\
    Lampe orange   & 4H 5L 5H \\
    Lampe grün     & 1H 1L 9H 3L \\
  \extracode
    \tablerules
    \vertlines[help lines]{0,1,...,14}
  \end{tikztimingtable}}
\end{zeitdiagramm}

\clearpage
\section{Fehlerfall 1: Not-Aus während der Förderung}
\label{sec:verlauf-notaus}

Während LÄUFT wird der Not-Aus an Modul~1 gedrückt, das Signal wechselt auf
low. Motor und Druckluft gehen sofort aus, die rote Lampe leuchtet. Nach dem
Entriegeln ist das Signal wieder high, und der Ablauf startet erneut über
BEREIT und VORLAUF.

\begin{zeitdiagramm}{Zeitverlauf bei Not-Aus an Modul~1}{fig:verlauf-notaus}
  \resizebox{\textwidth}{!}{%
  \begin{tikztimingtable}[timing/slope=0, xscale=2.6, yscale=1.4, font=\small]
    Zustand        & 1D{INIT} 1D{BER.} 2D{VORLAUF} 4D{LÄUFT} 3D{STOP} 1D{BER.} 2D{VORLAUF} \\
    Silo Bottom    & 14H \\
    Silo Top       & 14L \\
    Trichter       & 14L \\
    Not-Aus        & 8H 3L 3H \\
    M2 Silo voll   & 14L \\
    M2 Not-Aus     & 14H \\
    Motor          & 4L 4H 6L \\
    Druckluft      & 2L 6H 4L 2H \\
    Lampe rot      & 1H 7L 3H 3L \\
    Lampe orange   & 4H 7L 3H \\
    Lampe grün     & 1H 1L 6H 4L 2H \\
  \extracode
    \tablerules
    \vertlines[help lines]{0,1,...,14}
  \end{tikztimingtable}}
\end{zeitdiagramm}

\clearpage
\section{Fehlerfall 2: Trichter voll}
\label{sec:verlauf-trichter}

Der Trichter am Schlauch meldet \glqq voll\grqq, zum Beispiel weil der Schlauch
verstopft ist. Die Steuerung geht in STOP, bis der Trichter wieder frei ist.

\begin{zeitdiagramm}{Zeitverlauf bei vollem Trichter}{fig:verlauf-trichter}
  \resizebox{\textwidth}{!}{%
  \begin{tikztimingtable}[timing/slope=0, xscale=2.6, yscale=1.4, font=\small]
    Zustand        & 1D{INIT} 1D{BER.} 2D{VORLAUF} 4D{LÄUFT} 3D{STOP} 1D{BER.} 2D{VORLAUF} \\
    Silo Bottom    & 14H \\
    Silo Top       & 14L \\
    Trichter       & 8L 3H 3L \\
    Not-Aus        & 14H \\
    M2 Silo voll   & 14L \\
    M2 Not-Aus     & 14H \\
    Motor          & 4L 4H 6L \\
    Druckluft      & 2L 6H 4L 2H \\
    Lampe rot      & 1H 7L 3H 3L \\
    Lampe orange   & 4H 7L 3H \\
    Lampe grün     & 1H 1L 6H 4L 2H \\
  \extracode
    \tablerules
    \vertlines[help lines]{0,1,...,14}
  \end{tikztimingtable}}
\end{zeitdiagramm}

\clearpage
\section{Fehlerfall 3: Modul 2 meldet Silo voll}
\label{sec:verlauf-m2-voll}

Während LÄUFT meldet Modul~2 über das Handshake-Topic \texttt{M2\_Silo\_Full},
dass sein Silo voll ist. Im Gegensatz zu den Fehlerfällen~1 und~2 ist das kein
Sicherheitsfehler. Modul~1 geht deshalb nicht in STOP, sondern beendet die
Förderung geregelt: Band und Druckluft laufen 1~s nach, damit der Schlauch leer
geblasen wird. Danach wartet Modul~1 in BEREIT (orange). Sobald Modul~2 wieder
Platz meldet, startet der Ablauf neu über VORLAUF.

\begin{zeitdiagramm}{Zeitverlauf bei vollem Silo an Modul~2}{fig:verlauf-m2-voll}
  \resizebox{\textwidth}{!}{%
  \begin{tikztimingtable}[timing/slope=0, xscale=2.6, yscale=1.4, font=\small]
    Zustand        & 1D{INIT} 1D{BER.} 2D{VORLAUF} 2D{LÄUFT} 2D{NACHLAUF} 2D{BEREIT} 2D{VORLAUF} 2D{LÄUFT} \\
    Silo Bottom    & 14H \\
    Silo Top       & 14L \\
    Trichter       & 14L \\
    Not-Aus        & 14H \\
    M2 Silo voll   & 6L 4H 4L \\
    M2 Not-Aus     & 14H \\
    Motor          & 4L 4H 4L 2H \\
    Druckluft      & 2L 6H 2L 4H \\
    Lampe rot      & 1H 13L \\
    Lampe orange   & 4H 2L 6H 2L \\
    Lampe grün     & 1H 1L 6H 2L 4H \\
  \extracode
    \tablerules
    \vertlines[help lines]{0,1,...,14}
  \end{tikztimingtable}}
\end{zeitdiagramm}

\chapter{Kommunikation zu den Nachbarstationen}
\label{chap:kommunikation}

% Inhalt laut Vorgabe:
%  - Dokumentation der MQTT-Topics zur Kommunikation mit den Nachbarstationen
%  - inkl. JSON-Format

Modul~1 hat nur einen Nachbarn, und zwar Modul~2 (Storage), das den geförderten Mais
aufnimmt. Damit Modul~1 nicht in ein volles Lager fördert und bei einem Not-Aus an
Modul~2 stoppt, tauschen beide Module Signale über eigene Topics unterhalb von
\texttt{MV/VPS/DeliveryStorage/} aus. Diese Topics gehören nicht zur
Standardbelegung der Anlage, sondern wurden für die Kopplung angelegt.

\section{Topics}

\begin{table}[htbp]
  \centering
  \footnotesize
  \begin{tabularx}{\textwidth}{@{}l l l X@{}}
    \toprule
    \textbf{Topic} & \textbf{Richtung} & \textbf{Default} & \textbf{Bedeutung von \texttt{true}} \\
    \midrule
    \texttt{.../M2\_Silo\_Full}  & M2 $\rightarrow$ M1 & \texttt{false} & Silo von Modul~2 ist voll, Modul~1 darf nicht fördern \\
    \texttt{.../M2\_Emergency}   & M2 $\rightarrow$ M1 & \texttt{true}  & Not-Aus an Modul~2 ist \emph{nicht} gedrückt \\
    \bottomrule
  \end{tabularx}
  \caption[Handshake-Topics zwischen Modul~1 und Modul~2]{Handshake-Topics zwischen Modul~1 und Modul~2, Topic-Präfix \texttt{MV/VPS/DeliveryStorage/}}
  \label{tab:handshake}
\end{table}

Der Default gilt, solange Modul~1 noch keine Nachricht empfangen hat.
\texttt{M2\_Emergency} folgt derselben Logik wie der eigene Not-Aus
(Kapitel~\ref{sec:sensoren-aktoren}): Modul~2 leitet den Wert seines
Not-Aus-Sensors unverändert weiter, \texttt{true} heißt also \glqq kein Not-Aus\grqq.
Mit dem Default \texttt{true} fördert Modul~1 deshalb auch, solange von
Modul~2 noch keine Nachricht vorliegt. Beide Topics werdem mit \abkref{QoS}~0 publiziert.

\section{Sendezeitpunkt}

Die Kommunikation läuft nur in eine Richtung: Modul~2 sendet, Modul~1 liest. 
Handshakes von Modul~2 sind weitergeleitete Signale der eigenen Sensoren:

% raggedright: verhindert große Wortabstände wegen der langen Topic-Namen
\begin{itemize}[before=\raggedright]
  \item Trifft eine Nachricht von Silo Top an Modul~2 ein
        (\texttt{.../Sensors/Storage/Top}), sendet Modul~2 \texttt{M2\_Silo\_Full}.
  \item Trifft eine Nachricht vom Not-Aus an Modul~2 ein
        (\texttt{.../Sensors/Storage/Emergency}), sendet Modul~2 \texttt{M2\_Emergency}.
\end{itemize}

\section{Nachrichtenformat}

Die Handshake-Signale verwenden dasselbe \abkref{JSON}-Format wie die
Sensoren der Anlage: ein Objekt mit dem Wahrheitswert im Feld \texttt{value}
und dem Datentyp im Feld \texttt{dataType}. Modul~1 wertet nur \texttt{value}
aus, weitere Felder werden ignoriert.

\begin{lstlisting}[language=json, caption={Payload eines Handshake-Signals}, label={lst:handshake}]
{"value":true,"dataType":"Boolean"}
\end{lstlisting}

Zum Vergleich die Formate der eigenen Aktoren von Modul~1. Die Lampen erhalten
ein JSON-Objekt mit genau den drei Schlüsseln \texttt{Red}, \texttt{Orange}
und \texttt{Green}:

\begin{lstlisting}[language=json, caption={Payload der Lampen (Topic MV/VPS/Delivery/Lamps)}, label={lst:lampen}]
{"Red": false, "Orange": true, "Green": false}
\end{lstlisting}

Motor und Druckluft werden nicht mit JSON, sondern mit dem reinen Text
\texttt{true} bzw. \texttt{false} geschaltet (Topics \texttt{.../Transport/Actor}
und \texttt{.../Transport/Air}, Präfix \texttt{MV/VPS/Delivery/}). Alle Befehle von
Modul~1 werden ebenfalls mit QoS~0 gesendet.

\section{Abonnements von Modul 1}

\begin{itemize}
  \item \texttt{MV/VPS/Sensors/Delivery/\*} -- alle eigenen Sensoren
  \item \texttt{MV/VPS/DeliveryStorage/M2\_Silo\_Full}
  \item \texttt{MV/VPS/DeliveryStorage/M2\_Emergency}
\end{itemize}

\section{Publizierte Topics von Modul 1}

\begin{itemize}
  \item \texttt{MV/VPS/Delivery/Transport/Actor} -- Förderband (Motor)
  \item \texttt{MV/VPS/Delivery/Transport/Air} -- Druckluft
  \item \texttt{MV/VPS/Delivery/Lamps} -- Lampen rot, orange, grün
\end{itemize}

\chapter{Visuelle Signale}
\label{chap:signale}

% Inhalt laut Vorgabe:
%  - Dokumentation der visuellen Signale mit Bedeutung

Modul~1 zeigt seinen Zustand über eine Signalsäule mit drei Lampen an
(rot, orange, grün). Die Lampen werden beim Zustandswechsel über das Topic
\texttt{MV/VPS/Delivery/Lamps} gesetzt.

\begin{table}[htbp]
  \centering
  \renewcommand{\arraystretch}{1.5}
  \begin{tabularx}{\textwidth}{@{}l c c c X@{}}
    \toprule
    \textbf{Zustand} & \textbf{Rot} & \textbf{Orange} & \textbf{Grün} & \textbf{Bedeutung für den Bediener} \\
    \midrule
    INIT     & \lampe{lamperot}{1} & \lampe{lampeorange}{1} & \lampe{lampegruen}{1} & Lampentest beim Start \\
    BEREIT   & \lampe{lamperot}{0} & \lampe{lampeorange}{1} & \lampe{lampegruen}{0} & Wartet auf Material im Silo bzw. Freigabe von Modul~2 \\
    VORLAUF  & \lampe{lamperot}{0} & \lampe{lampeorange}{1} & \lampe{lampegruen}{1} & Druckluft läuft vor, Förderband startet gleich \\
    LÄUFT    & \lampe{lamperot}{0} & \lampe{lampeorange}{0} & \lampe{lampegruen}{1} & Normale Förderung \\
    NACHLAUF & \lampe{lamperot}{0} & \lampe{lampeorange}{1} & \lampe{lampegruen}{1} & Förderung endet, Band und Luft laufen aus \\
    STOP     & \lampe{lamperot}{1} & \lampe{lampeorange}{0} & \lampe{lampegruen}{0} & Sicherheitsstopp: Not-Aus (M1 oder M2) oder Trichter voll \\
    \bottomrule
  \end{tabularx}
  \caption[Lampenzustände von Modul~1 und ihre Bedeutung]{Lampenzustände von Modul~1 und ihre Bedeutung
    (\protect\lampe{gray}{1}~an, \protect\lampe{gray}{0}~aus)}
  \label{tab:lampen}
\end{table}

VORLAUF und NACHLAUF zeigen dasselbe Lampenbild, weil in beiden Fällen
Druckluft und gegebenenfalls das Band laufen, aber keine reguläre Förderung
stattfindet.


\clearpage
\phantomsection
\addcontentsline{toc}{chapter}{Literaturverzeichnis}
\markboth{Literaturverzeichnis}{}
\begin{thebibliography}{9}
  \raggedright
  \bibitem{doku-popcorn}
  inIT / TH OWL: Maschinennahe Vernetzung Praktikum -- von MQTT bis zu Node-Red
  (\texttt{popcornUA\_doku\_studis.pdf}). Praktikumsunterlage, \Semester.

  \bibitem{vps-simple-view}
  inIT / TH OWL: VPS Simple View -- Übersicht der MQTT-Topics aller Module.
  Online: \url{http://frankfort.init.local:1337/simple.php} (nur im Hochschulnetz),
  Zugriff am 15.10.2026.

  \bibitem{paho}
  Eclipse Foundation: Eclipse Paho MQTT Python Client, Version 2.1.
  Online: \url{https://eclipse.dev/paho/files/paho.mqtt.python/html/}, Zugriff am 13.11.2026.

  \bibitem{mqtt311}
  OASIS: MQTT Version 3.1.1. OASIS Standard, 2014.
  Online: \url{https://docs.oasis-open.org/mqtt/mqtt/v3.1.1/mqtt-v3.1.1.html}, Zugriff am 13.11.2026.

  \bibitem{rfc8259}
  Bray, T.: RFC 8259, The JavaScript Object Notation (JSON) Data Interchange Format.
  IETF, 2017. Online: \url{https://www.rfc-editor.org/rfc/rfc8259}, Zugriff am 13.11.2026.
\end{thebibliography}

\end{document}
